\chapter{Superconductivity and circuit quantisation}
\label{ch:circuit-quantisation}

Everything in Parts~I--III was platform-independent: a qubit was an
abstract two-level system, a cavity an abstract harmonic oscillator,
and the Jaynes--Cummings coupling a number $g$ we never had to
derive from hardware. Part~IV makes the platform concrete. The
qubits and resonators of a superconducting quantum processor are
\emph{electrical circuits} --- capacitors, inductors, and Josephson
junctions --- cooled until a macroscopic number of electrons condense
into a single coherent quantum state and the circuit's collective
coordinate, the flux through an inductor or the charge on a
capacitor, behaves quantum-mechanically. To build that picture we
need two things this chapter supplies. First, \emph{why} a metal
becomes a dissipationless quantum conductor at all: the
Bardeen--Cooper--Schrieffer (BCS) theory of pairing, which produces
the energy gap $2\Delta$ that Part~VI has been invoking inline ever
since \cref{ch:qp-detection}. Second, \emph{how} to quantise a
circuit: the canonical procedure that promotes node flux and charge
to conjugate operators and turns an $LC$ resonator into the quantum
harmonic oscillator of \cref{stmt:qho-spectrum}. The chapter closes
by framing what a superconducting platform must still deliver --- the
DiVincenzo criteria --- and identifying the one ingredient the linear
$LC$ circuit lacks: anharmonicity, the Josephson junction of
\cref{ch:josephson}.

% --------------------------------------------------------------
\section{Superconductivity as a macroscopic quantum phenomenon}
\label{sec:superconductivity-phenom}

\begin{phenomenon}
Below a critical temperature $T_c$ many metals lose all DC
electrical resistance: a current set up in a superconducting ring
persists for years without measurable decay. They also expel
magnetic flux from their interior (the Meissner effect), so a
superconductor is not merely a perfect conductor but a distinct
thermodynamic phase. Two further clues point to the microscopic
mechanism. The critical temperature depends on the ionic mass
($T_c\propto M^{-1/2}$, the \emph{isotope effect}), implicating the
lattice vibrations; and magnetic flux trapped through a
superconducting ring comes only in integer multiples of
$\Phi_0=h/2e$, half the value $h/e$ one would expect for single
electrons. The factor of two in the charge is the signature of
\emph{pairing}.
\end{phenomenon}

\begin{statement}[\textbf{Flux quantisation and the macroscopic wavefunction}]
\label{stmt:flux-quantisation}
A superconductor is described by a single complex \emph{order
parameter} (macroscopic wavefunction)
\begin{equation}
\label{eq:order-parameter}
\psi(\mathbf r) \;=\; |\psi(\mathbf r)|\,e^{i\varphi(\mathbf r)},
\end{equation}
whose modulus sets the density of the condensed charge carriers and
whose phase $\varphi$ is rigid across the sample --- a single quantum
state occupied by a macroscopic number of carriers. Because the
carriers have charge $2e$ (Cooper pairs), single-valuedness of
$\psi$ around a closed loop forces the enclosed flux to be quantised,
\begin{equation}
\label{eq:flux-quantum}
\oint \nabla\varphi\cdot\dd\mathbf l = 2\pi n
\quad\Longrightarrow\quad
\Phi = n\,\Phi_0,
\qquad
\Phi_0 \equiv \frac{h}{2e}\approx 2.07\times10^{-15}\,\text{Wb}.
\end{equation}
The phase $\varphi$ is the dynamical variable of the superconducting
condensate; it is the same phase that will appear across a Josephson
junction in \cref{ch:josephson} and that, together with the charge,
becomes a quantum operator later in this chapter.
\end{statement}

\paragraph{Why the phase matters.}
The rigidity of $\varphi$ is what makes a superconducting circuit a
clean quantum system. In a normal metal the electronic degrees of
freedom are a vast incoherent bath; in a superconductor they are
locked into one coherent state, and the only low-energy dynamics left
is that of the collective phase and the charge conjugate to it. A
circuit built from superconducting elements therefore has just a
handful of relevant quantum coordinates, not Avogadro's number of
them --- which is precisely what allows an $LC$ loop to behave as a
single harmonic oscillator and a Josephson circuit as a single qubit.

\paragraph{Transition.}
The phenomenology fixes \emph{what} a superconductor is --- a
macroscopic quantum state of charge-$2e$ carriers with a rigid phase
--- but not \emph{why} electrons, which repel one another, should
pair at all. That is the content of BCS theory, to which we now turn.

% --------------------------------------------------------------
\section{The Cooper instability and the BCS ground state}
\label{sec:cooper-bcs}

\begin{phenomenon}
Electrons repel through the Coulomb interaction, yet they pair. The
resolution is that the relevant interaction near the Fermi surface is
not the bare Coulomb force but an \emph{effective} one mediated by
the lattice: one electron distorts the positive-ion lattice, and a
second electron is attracted to the transient excess of positive
charge left behind. The attraction is retarded and weak, but Cooper
showed that \emph{any} attractive interaction, however weak,
destabilises the normal Fermi sea: a single pair of electrons added
just above a filled Fermi sea forms a bound state. The ground state
must therefore reorganise into a coherent superposition of such
pairs.
\end{phenomenon}

\begin{statement}[\textbf{Reduced BCS pairing Hamiltonian}~\cite{bcs1957}]
\label{stmt:bcs-hamiltonian}
Keeping only the interaction between electrons of opposite momentum
and spin --- the channel Cooper found unstable --- the reduced BCS
Hamiltonian is
\begin{equation}
\label{eq:bcs-hamiltonian}
\op H_{\rm BCS}
= \sum_{\mathbf k,\sigma}\xi_{\mathbf k}\,
  \hat c^\dagger_{\mathbf k\sigma}\hat c_{\mathbf k\sigma}
\;-\; V\!\!\sum_{\mathbf k,\mathbf k'}
  \hat c^\dagger_{\mathbf k'\uparrow}\hat c^\dagger_{-\mathbf k'\downarrow}
  \hat c_{-\mathbf k\downarrow}\hat c_{\mathbf k\uparrow},
\end{equation}
where $\hat c^\dagger_{\mathbf k\sigma}$ creates an electron of
momentum $\mathbf k$ and spin $\sigma$, $\xi_{\mathbf k}=
\varepsilon_{\mathbf k}-\mu$ is the single-particle energy measured
from the chemical potential $\mu$, and $V>0$ is the strength of the
attractive interaction, taken constant for
$|\xi_{\mathbf k}|<\hbar\omega_D$ (a thin shell of width set by the
Debye frequency $\omega_D$) and zero otherwise. The operators obey
the fermionic anticommutation relations
$\{\hat c_{\mathbf k\sigma},\hat c^\dagger_{\mathbf k'\sigma'}\}
=\delta_{\mathbf k\mathbf k'}\delta_{\sigma\sigma'}$.
\end{statement}

\begin{statement}[\textbf{BCS variational ground state}]
\label{stmt:bcs-ground-state}
The ground state of \cref{eq:bcs-hamiltonian} is the coherent pair
condensate
\begin{equation}
\label{eq:bcs-state}
\ket{\rm BCS}
= \prod_{\mathbf k}
  \bigl(u_{\mathbf k}
   + v_{\mathbf k}\,\hat c^\dagger_{\mathbf k\uparrow}
     \hat c^\dagger_{-\mathbf k\downarrow}\bigr)\ket0,
\qquad
|u_{\mathbf k}|^2 + |v_{\mathbf k}|^2 = 1,
\end{equation}
in which each pair state $(\mathbf k\uparrow,-\mathbf k\downarrow)$ is
empty with amplitude $u_{\mathbf k}$ and occupied with amplitude
$v_{\mathbf k}$. Unlike a filled Fermi sea, $\ket{\rm BCS}$ is not an
eigenstate of particle number: it is a superposition of states with
different numbers of pairs, with a well-defined relative \emph{phase}
between them --- the microscopic origin of the order parameter
\cref{eq:order-parameter}.
\end{statement}

\paragraph{Derivation (mean-field decoupling).}
The quartic interaction in \cref{eq:bcs-hamiltonian} is intractable
exactly, but the coherence of $\ket{\rm BCS}$ makes the pair operator
$\hat c_{-\mathbf k\downarrow}\hat c_{\mathbf k\uparrow}$ acquire a
nonzero expectation value. Define the \emph{gap parameter}
\begin{equation}
\label{eq:gap-definition}
\Delta \;\equiv\; V\sum_{\mathbf k}
  \expval{\hat c_{-\mathbf k\downarrow}\hat c_{\mathbf k\uparrow}}.
\end{equation}
Writing each pair operator as its mean plus a fluctuation,
$\hat c_{-\mathbf k\downarrow}\hat c_{\mathbf k\uparrow}
=\expval{\cdots}+\delta$, and dropping the (small) product of two
fluctuations, the interaction linearises into the mean-field
Hamiltonian
\begin{equation}
\label{eq:bcs-meanfield}
\op H_{\rm MF}
= \sum_{\mathbf k,\sigma}\xi_{\mathbf k}\,
  \hat c^\dagger_{\mathbf k\sigma}\hat c_{\mathbf k\sigma}
  \;-\;\sum_{\mathbf k}\Bigl(
  \Delta\,\hat c^\dagger_{\mathbf k\uparrow}
    \hat c^\dagger_{-\mathbf k\downarrow} + \hc\Bigr)
  \;+\;\frac{|\Delta|^2}{V}.
\end{equation}
This is quadratic and can be diagonalised exactly --- the subject of
the next section --- and the diagonalisation feeds back through
\cref{eq:gap-definition} to fix $\Delta$ self-consistently.

\paragraph{Transition.}
The mean-field Hamiltonian \cref{eq:bcs-meanfield} mixes creation and
annihilation operators ($\hat c^\dagger\hat c^\dagger$ terms), so the
electron is no longer the natural excitation. The Bogoliubov
transformation finds the operators that \emph{do} diagonalise it, and
in doing so produces both the quasiparticle spectrum and the
self-consistent equation for the gap.

% --------------------------------------------------------------
\section{The gap equation and Bogoliubov quasiparticles}
\label{sec:gap-bogoliubov}

\begin{statement}[\textbf{Bogoliubov quasiparticles and the gap equation}]
\label{stmt:bogoliubov}
The mean-field Hamiltonian \cref{eq:bcs-meanfield} is diagonalised by
the \emph{Bogoliubov transformation} to new fermionic operators
\begin{equation}
\label{eq:bogoliubov-transform}
\hat\gamma_{\mathbf k0}
  = u_{\mathbf k}\hat c_{\mathbf k\uparrow}
  - v_{\mathbf k}\hat c^\dagger_{-\mathbf k\downarrow},
\qquad
\hat\gamma_{\mathbf k1}
  = u_{\mathbf k}\hat c_{-\mathbf k\downarrow}
  + v_{\mathbf k}\hat c^\dagger_{\mathbf k\uparrow},
\end{equation}
with the coherence factors
\begin{equation}
\label{eq:coherence-factors}
u_{\mathbf k}^2 = \frac12\Bigl(1+\frac{\xi_{\mathbf k}}{E_{\mathbf k}}\Bigr),
\qquad
v_{\mathbf k}^2 = \frac12\Bigl(1-\frac{\xi_{\mathbf k}}{E_{\mathbf k}}\Bigr),
\qquad
E_{\mathbf k}=\sqrt{\xi_{\mathbf k}^2+\Delta^2}.
\end{equation}
In terms of these the Hamiltonian becomes a free gas of
\emph{quasiparticles},
\begin{equation}
\label{eq:bcs-diagonal}
\op H_{\rm MF}
= E_0 + \sum_{\mathbf k}E_{\mathbf k}\bigl(
  \hat\gamma^\dagger_{\mathbf k0}\hat\gamma_{\mathbf k0}
  + \hat\gamma^\dagger_{\mathbf k1}\hat\gamma_{\mathbf k1}\bigr),
\end{equation}
whose excitation energy $E_{\mathbf k}=\sqrt{\xi_{\mathbf k}^2+\Delta^2}$
never falls below $\Delta$. The minimum energy to create
\emph{two} quasiparticles --- to break one Cooper pair --- is
therefore $2\Delta$. Self-consistency \cref{eq:gap-definition} closes
into the \emph{gap equation}
\begin{equation}
\label{eq:gap-equation}
1 = V\sum_{\mathbf k}\frac{1}{2E_{\mathbf k}}
  = N(0)V\!\int_0^{\hbar\omega_D}\!\frac{\dd\xi}{\sqrt{\xi^2+\Delta^2}}
  = N(0)V\,\operatorname{arcsinh}\!\frac{\hbar\omega_D}{\Delta},
\end{equation}
with $N(0)$ the normal-state density of states at the Fermi level
(per spin). In the weak-coupling limit $N(0)V\ll1$ this gives the
celebrated BCS gap
\begin{equation}
\label{eq:bcs-gap}
\Delta \;=\; \frac{\hbar\omega_D}{\sinh\!\bigl(1/N(0)V\bigr)}
  \;\approx\; 2\hbar\omega_D\,e^{-1/N(0)V}.
\end{equation}
\end{statement}

\paragraph{Derivation (diagonalisation).}
The transformation \cref{eq:bogoliubov-transform} is canonical --- it
preserves the fermionic anticommutators --- provided
$u_{\mathbf k}^2+v_{\mathbf k}^2=1$, which is why the same constraint
appeared in the variational state \cref{eq:bcs-state}. Substituting
the inverse of \cref{eq:bogoliubov-transform} into
\cref{eq:bcs-meanfield} produces, besides the diagonal
$\hat\gamma^\dagger\hat\gamma$ terms, off-diagonal pieces
$\propto\hat\gamma^\dagger_{\mathbf k0}\hat\gamma^\dagger_{\mathbf k1}$.
Demanding that these vanish fixes the mixing angle: with
$u_{\mathbf k}=\cos\theta_{\mathbf k}$, $v_{\mathbf k}=\sin\theta_{\mathbf k}$,
the off-diagonal terms cancel when
$\tan2\theta_{\mathbf k}=\Delta/\xi_{\mathbf k}$, equivalent to
\cref{eq:coherence-factors} and to
$2u_{\mathbf k}v_{\mathbf k}=\Delta/E_{\mathbf k}$. The surviving
diagonal coefficient is exactly $E_{\mathbf k}$. Feeding
$\expval{\hat c_{-\mathbf k\downarrow}\hat c_{\mathbf k\uparrow}}
= u_{\mathbf k}v_{\mathbf k}=\Delta/2E_{\mathbf k}$ (at zero
temperature, where the quasiparticle vacuum is empty) back into the
definition \cref{eq:gap-definition} gives the gap equation
\cref{eq:gap-equation}; the factor $\Delta$ cancels on both sides,
leaving the condition on $V$ and $\omega_D$ alone.

\begin{remark}
The gap \cref{eq:bcs-gap} is non-analytic in the coupling $V$: the
function $e^{-1/N(0)V}$ has an essential singularity at $V=0$, so no
order of perturbation theory in $V$ could ever produce it. This is
why the normal Fermi sea is unstable to \emph{any} attraction, and
why BCS required a non-perturbative variational ansatz rather than a
diagrammatic expansion around the free gas.
\end{remark}

\begin{application}[The gap as the energy scale of Part~VI]
\label{app:gap-partvi}
The quantity $2\Delta$ derived here is the same gap that governs the
quasiparticle detectors of Part~VI. In aluminium $2\Delta\approx
0.34\,\text{meV}$; an athermal phonon with energy above $2\Delta$ can
break a Cooper pair and inject the quasiparticle excitations that a
transition-edge sensor or kinetic-inductance detector reads out
(\cref{ch:qp-detection}). The exponential smallness of $\Delta$
relative to the Debye energy --- a few kelvin against a few hundred
--- is what makes these detectors sensitive to milli-electronvolt
energy depositions, far below any normal-metal calorimeter. The same
gap also sets the ceiling on a superconducting qubit's operating
frequency: drive photons must stay well below $2\Delta/\hbar$ or they
break pairs and poison coherence, the quasiparticle-poisoning channel
of \cref{ch:qp-detection}.
\end{application}

\paragraph{Transition.}
BCS theory explains why a superconducting circuit is dissipationless
at low frequency: there are no electronic excitations below $2\Delta$
to carry loss, so currents flow without resistance and the only
surviving low-energy dynamics is that of the collective phase and
charge. We now treat those collective coordinates as the dynamical
variables of an electrical circuit and quantise them.

% --------------------------------------------------------------
\section{Lumped-element circuit quantisation}
\label{sec:circuit-quantisation}

\begin{phenomenon}
A capacitor and an inductor wired in a loop exchange energy back and
forth --- charge sloshes onto the capacitor, discharges through the
inductor, and recharges with opposite sign --- exactly as a mass on a
spring exchanges kinetic and potential energy. Classically this $LC$
circuit rings at $\omega=1/\sqrt{LC}$ and, in a normal metal, damps
out through resistance. Built from superconductor and cooled below
the point where $\hbar\omega\gg k_BT$, the circuit instead becomes a
nearly lossless oscillator whose single collective coordinate must be
described quantum-mechanically: its energy comes in discrete quanta,
and its flux and charge cannot be sharp simultaneously.
\end{phenomenon}

\begin{statement}[\textbf{Canonical quantisation of an $LC$ circuit}~\cite{devoret1995,vool2017}]
\label{stmt:circuit-quantisation}
Take as coordinate the \emph{node flux}
$\Phi(t)=\int_{-\infty}^{t}V(t')\dd t'$, the time integral of the
voltage across the inductor. The capacitor stores energy
$Q^2/2C$ and the inductor $\Phi^2/2L$, so the circuit Lagrangian is
\begin{equation}
\label{eq:lc-lagrangian}
\mathcal L = \underbrace{\tfrac12 C\dot\Phi^2}_{\text{``kinetic''}}
  - \underbrace{\frac{\Phi^2}{2L}}_{\text{``potential''}},
\end{equation}
with $\dot\Phi=V$. The momentum conjugate to $\Phi$ is the charge
$Q=\partial\mathcal L/\partial\dot\Phi = C\dot\Phi$, and a Legendre
transform gives the Hamiltonian
\begin{equation}
\label{eq:lc-hamiltonian}
H = Q\dot\Phi - \mathcal L = \frac{Q^2}{2C} + \frac{\Phi^2}{2L}.
\end{equation}
Promoting the conjugate pair to operators with
\begin{equation}
\label{eq:flux-charge-commutator}
\comm{\op\Phi}{\op Q} = i\hbar
\end{equation}
quantises the circuit. The mapping to a mechanical oscillator is
exact: flux $\Phi\leftrightarrow$ position, charge
$Q\leftrightarrow$ momentum, capacitance $C\leftrightarrow$ mass,
inverse inductance $1/L\leftrightarrow$ spring constant.
\end{statement}

\paragraph{Derivation.}
Kirchhoff's laws for the loop read $I_C=I_L$ (current conservation)
and the same voltage $V$ across both elements. With $\Phi=\int V\dd t$
the inductor current is $I_L=\Phi/L$ (from $V=L\dot I_L$) and the
capacitor current is $I_C=C\dot V=C\ddot\Phi$. Equating them gives
$C\ddot\Phi+\Phi/L=0$, the harmonic equation of motion with
$\omega=1/\sqrt{LC}$. This same equation follows as the
Euler--Lagrange equation of \cref{eq:lc-lagrangian},
$\frac{\dd}{\dd t}\frac{\partial\mathcal L}{\partial\dot\Phi}
-\frac{\partial\mathcal L}{\partial\Phi}=0$, confirming the
Lagrangian is correct. The conjugate momentum
$Q=\partial\mathcal L/\partial\dot\Phi=C\dot\Phi$ is just the charge
$Q=CV$ on the capacitor, so the canonical pair $(\Phi,Q)$ is
(flux, charge) --- physical, measurable quantities, not abstract
coordinates.

\begin{remark}
The choice of node flux as ``position'' and charge as ``momentum'' is
a convention, not a law: one may equally take charge as coordinate
and flux as momentum (the \emph{charge representation}, natural for
the Cooper-pair box of \cref{ch:transmon}). The two are related by
the exchange $\Phi\leftrightarrow Q$, $L\leftrightarrow C$ ---
electrical \emph{duality} --- and \cref{eq:flux-charge-commutator} is
symmetric under it up to a sign. Which representation is convenient
depends on whether the dominant nonlinearity is inductive (a
Josephson junction, flux-like) or capacitive (a small island,
charge-like).
\end{remark}

\paragraph{Transition.}
With the circuit cast as a canonical Hamiltonian system, the same
ladder-operator machinery that solved the mechanical oscillator in
\cref{ch:model-systems} and quantised the electromagnetic field in
\cref{ch:em-quantisation} applies verbatim. We introduce it and read
off the quantum properties of the resonator.

% --------------------------------------------------------------
\section{The \texorpdfstring{$LC$}{LC} oscillator as a quantum harmonic oscillator}
\label{sec:lc-oscillator}

\begin{statement}[\textbf{Quantised $LC$ resonator}]
\label{stmt:lc-oscillator}
Define the characteristic impedance $Z=\sqrt{L/C}$ and the
zero-point fluctuations of flux and charge,
\begin{equation}
\label{eq:zpf}
\Phi_{\rm zpf}=\sqrt{\frac{\hbar Z}{2}},
\qquad
Q_{\rm zpf}=\sqrt{\frac{\hbar}{2Z}},
\qquad
\Phi_{\rm zpf}Q_{\rm zpf}=\frac{\hbar}{2}.
\end{equation}
Introducing ladder operators through
\begin{equation}
\label{eq:lc-ladder}
\op\Phi=\Phi_{\rm zpf}(\hat a+\hat a^\dagger),
\qquad
\op Q=-iQ_{\rm zpf}(\hat a-\hat a^\dagger),
\qquad
\comm{\hat a}{\hat a^\dagger}=\id,
\end{equation}
the Hamiltonian \cref{eq:lc-hamiltonian} becomes
\begin{equation}
\label{eq:lc-qho}
\op H = \hbar\omega\Bigl(\hat a^\dagger\hat a + \tfrac12\Bigr),
\qquad
\omega=\frac{1}{\sqrt{LC}},
\end{equation}
the quantum harmonic oscillator of \cref{stmt:qho-spectrum}. The
resonator's energy is quantised in units of $\hbar\omega$; one quantum
is a single microwave photon, and the vacuum carries irreducible
flux and charge fluctuations $\Phi_{\rm zpf}$, $Q_{\rm zpf}$ obeying
the uncertainty bound $\Delta\Phi\,\Delta Q\geq\hbar/2$.
\end{statement}

\paragraph{The same oscillator, three times over.}
\Cref{eq:lc-qho} is identical in form to the mechanical oscillator of
\cref{stmt:qho-spectrum} and to a single mode of the quantised
electromagnetic field in \cref{stmt:field-oscillators}: the
quadratures of \cref{stmt:quadratures} are here the flux and charge,
and the vacuum noise of \cref{stmt:quadratures} is the zero-point
current and voltage noise of the resonator. A superconducting $LC$
circuit is therefore a piece of cavity QED on a chip --- the
``cavity'' of \cref{ch:jaynes-cummings,ch:dispersive} realised as a
lithographically patterned resonator rather than a pair of mirrors.
This is what lets the abstract coupling $g$ of the Jaynes--Cummings
model acquire a concrete value in \cref{ch:circuit-qed}, computed from
circuit capacitances.

\paragraph{Typical numbers.}
For a resonator at $\omega/2\pi=5\,$GHz with a characteristic
impedance $Z\approx50\,\Omega$ (the standard microwave value), the
zero-point flux is $\Phi_{\rm zpf}=\sqrt{\hbar Z/2}\sim10^{-16}\,$Wb,
a small fraction of the flux quantum $\Phi_0$, while
$Q_{\rm zpf}=\sqrt{\hbar/2Z}$ is a fraction of a single electron
charge. The quantum of energy $\hbar\omega\approx20\,\mu$eV
corresponds to a temperature $\hbar\omega/k_B\approx240\,$mK, so a
dilution refrigerator at $10$--$20\,$mK holds the resonator firmly in
its ground state --- the precondition for the vacuum-state cavity
QED of Part~III to apply.

\paragraph{Why a linear resonator is not yet a qubit.}
The spectrum \cref{eq:lc-qho} is a uniform ladder: every transition
$\ket n\to\ket{n+1}$ has the identical energy $\hbar\omega$. A drive
tuned to $\ket0\to\ket1$ is therefore equally resonant with
$\ket1\to\ket2$, $\ket2\to\ket3$, and so on, so it cannot confine the
dynamics to the lowest two levels: the linear $LC$ circuit is a
perfect oscillator but a useless qubit. Isolating a two-level
subspace requires \emph{anharmonicity} --- an unequal level spacing
--- which a linear circuit element cannot provide. The unique
nonlinear, non-dissipative circuit element is the Josephson junction,
the subject of \cref{ch:josephson}; supplying the anharmonicity that
turns the resonator of this chapter into the transmon qubit of
\cref{ch:transmon} is its central role.

\paragraph{The DiVincenzo criteria.}
Stepping back, what must any platform provide to be a quantum
computer? DiVincenzo's checklist is the standard framing:
(i)~well-characterised qubits in a scalable architecture;
(ii)~the ability to initialise to a known state;
(iii)~long coherence times relative to gate times;
(iv)~a universal set of gates; and
(v)~qubit-specific measurement. Superconducting circuits address each
in turn through the rest of Part~IV: the transmon supplies the qubit
(\cref{ch:transmon}); thermalisation in a cold cavity supplies
initialisation; the long $T_1,T_2$ of Part~II, engineered in hardware
(\cref{ch:experimental-basics}), supply coherence; resonant drives and
tunable couplers supply single- and two-qubit gates
(\cref{ch:circuit-qed}); and the dispersive readout of
\cref{ch:dispersive} supplies measurement. The logical layer built on
top of these primitives --- gates, algorithms, and error correction
--- is the subject of Part~V (\cref{ch:quantum-information-processing}).

% --------------------------------------------------------------
This chapter laid the two foundations of the superconducting
platform. BCS theory showed that an arbitrarily weak phonon-mediated
attraction condenses electrons into Cooper pairs, opening a gap
$2\Delta$ in the excitation spectrum; the Bogoliubov transformation
gave the quasiparticles and the self-consistent gap equation, and
fixed the energy scale that Part~VI's detectors exploit and that
limits qubit coherence. Canonical quantisation then turned a
superconducting $LC$ loop into a quantum harmonic oscillator, with
node flux and charge as conjugate operators and microwave photons as
its quanta.

What the linear resonator lacks is anharmonicity, and the next
chapter supplies it. The Josephson junction --- a weak link between
two superconductors across which the condensate phase of
\cref{stmt:flux-quantisation} can tunnel --- is the nonlinear,
dissipationless inductor that breaks the uniform $LC$ ladder and makes
a circuit into a qubit.

\clearpage
\begin{chaptersummary}

\paragraph{Superconductivity.}
A superconductor is a macroscopic quantum state of charge-$2e$ Cooper
pairs with a rigid phase $\varphi$; flux is quantised in units of
$\Phi_0=h/2e$.

\paragraph{BCS pairing.}
The reduced Hamiltonian
$\op H_{\rm BCS}=\sum_{\mathbf k\sigma}\xi_{\mathbf k}\hat c^\dagger\hat c
-V\sum\hat c^\dagger_{\mathbf k'\uparrow}\hat c^\dagger_{-\mathbf k'\downarrow}
\hat c_{-\mathbf k\downarrow}\hat c_{\mathbf k\uparrow}$
has the coherent ground state
$\ket{\rm BCS}=\prod_{\mathbf k}(u_{\mathbf k}+v_{\mathbf k}
\hat c^\dagger_{\mathbf k\uparrow}\hat c^\dagger_{-\mathbf k\downarrow})\ket0$.

\paragraph{Quasiparticles and the gap.}
The Bogoliubov transformation gives quasiparticles of energy
$E_{\mathbf k}=\sqrt{\xi_{\mathbf k}^2+\Delta^2}\geq\Delta$, with
coherence factors
$u_{\mathbf k}^2,v_{\mathbf k}^2=\tfrac12(1\pm\xi_{\mathbf k}/E_{\mathbf k})$.
The gap equation $1=N(0)V\operatorname{arcsinh}(\hbar\omega_D/\Delta)$
yields $\Delta\approx2\hbar\omega_D e^{-1/N(0)V}$; breaking a pair
costs $2\Delta$ (e.g.\ $0.34\,$meV in Al).

\paragraph{Circuit quantisation.}
Node flux $\Phi=\int V\dd t$ and charge $Q=C\dot\Phi$ are conjugate,
$\comm{\op\Phi}{\op Q}=i\hbar$; an $LC$ loop has
$\op H=Q^2/2C+\Phi^2/2L$, with $C\leftrightarrow$ mass,
$1/L\leftrightarrow$ spring.

\paragraph{$LC$ as a quantum oscillator.}
With $Z=\sqrt{L/C}$, $\Phi_{\rm zpf}=\sqrt{\hbar Z/2}$,
$Q_{\rm zpf}=\sqrt{\hbar/2Z}$, the circuit is
$\op H=\hbar\omega(\hat a^\dagger\hat a+\tfrac12)$,
$\omega=1/\sqrt{LC}$ --- a microwave-photon oscillator. Its uniform
ladder cannot host a qubit; anharmonicity from a Josephson junction
(\cref{ch:josephson}) is required.

\end{chaptersummary}

% --------------------------------------------------------------
\section*{Exercises}
\addcontentsline{toc}{section}{Exercises}

\begin{exercise}[The flux quantum]
Evaluate $\Phi_0=h/2e$ numerically and explain why the charge in the
denominator is $2e$ rather than $e$. What would change if the
condensate carriers were single electrons?
\begin{solution}
$\Phi_0=h/2e=(6.626\times10^{-34})/(2\cdot1.602\times10^{-19})
=2.07\times10^{-15}\,$Wb. The $2e$ comes from the Cooper-pair charge:
the macroscopic wavefunction \cref{eq:order-parameter} describes
\emph{pairs}, and single-valuedness of its phase around a loop
quantises the flux in units of $h/(\text{carrier charge})$. Single
electrons would give $h/e=2\Phi_0$; the experimentally observed value
$h/2e$ was an early confirmation of pairing.
\end{solution}
\end{exercise}

\begin{exercise}[Coherence factors and normalisation]
Verify that $u_{\mathbf k}^2+v_{\mathbf k}^2=1$ for the coherence
factors \cref{eq:coherence-factors}, and compute the limits of
$v_{\mathbf k}^2$ far below ($\xi_{\mathbf k}\to-\infty$) and far above
($\xi_{\mathbf k}\to+\infty$) the Fermi surface. Interpret.
\begin{solution}
$u_{\mathbf k}^2+v_{\mathbf k}^2
=\tfrac12(1+\xi/E)+\tfrac12(1-\xi/E)=1$. As
$\xi\to-\infty$ (deep below $E_F$), $\xi/E\to-1$ so
$v_{\mathbf k}^2\to1$: states far below the Fermi level are fully
occupied, as in the normal metal. As $\xi\to+\infty$,
$v_{\mathbf k}^2\to0$: states far above are empty. The pairing only
reshuffles occupation within a shell of width $\sim\Delta$ around
$E_F$, where $v_{\mathbf k}^2$ smoothly interpolates from $1$ to $0$
(passing through $\tfrac12$ at $\xi=0$).
\end{solution}
\end{exercise}

\begin{exercise}[The pair amplitude is peaked at the Fermi surface]
Show that the pair amplitude $u_{\mathbf k}v_{\mathbf k}$ equals
$\Delta/2E_{\mathbf k}$ and is maximal at $\xi_{\mathbf k}=0$. Why does
this mean pairing is a Fermi-surface phenomenon?
\begin{solution}
From \cref{eq:coherence-factors},
$u_{\mathbf k}v_{\mathbf k}=\sqrt{u_{\mathbf k}^2 v_{\mathbf k}^2}
=\sqrt{\tfrac14(1-\xi^2/E^2)}=\tfrac12\sqrt{(E^2-\xi^2)/E^2}
=\tfrac12\sqrt{\Delta^2/E^2}=\Delta/2E_{\mathbf k}$. This is largest
when $E_{\mathbf k}$ is smallest, i.e.\ at $\xi_{\mathbf k}=0$ where
$E_{\mathbf k}=\Delta$. Since $u_{\mathbf k}v_{\mathbf k}$ is the
amplitude for a pair to scatter in/out of $(\mathbf k,-\mathbf k)$, the
condensate is built predominantly from states within $\sim\Delta$ of
the Fermi surface --- consistent with the thin interaction shell
$|\xi|<\hbar\omega_D$ in \cref{eq:bcs-hamiltonian}.
\end{solution}
\end{exercise}

\begin{exercise}[Solving the gap equation]
Carry out the integral in \cref{eq:gap-equation} to obtain the
$\operatorname{arcsinh}$ form, then take the weak-coupling limit to
recover \cref{eq:bcs-gap}.
\begin{solution}
$\int_0^{\hbar\omega_D}\dd\xi/\sqrt{\xi^2+\Delta^2}
=\operatorname{arcsinh}(\xi/\Delta)\big|_0^{\hbar\omega_D}
=\operatorname{arcsinh}(\hbar\omega_D/\Delta)$, so
$1=N(0)V\operatorname{arcsinh}(\hbar\omega_D/\Delta)$. Inverting,
$\operatorname{arcsinh}(\hbar\omega_D/\Delta)=1/N(0)V$, i.e.\
$\hbar\omega_D/\Delta=\sinh(1/N(0)V)$, giving
$\Delta=\hbar\omega_D/\sinh(1/N(0)V)$. For $N(0)V\ll1$ the argument is
large and $\sinh x\approx\tfrac12 e^{x}$, so
$\Delta\approx2\hbar\omega_D e^{-1/N(0)V}$.
\end{solution}
\end{exercise}

\begin{exercise}[Non-analyticity of the gap]
Show that $\Delta(V)=2\hbar\omega_D e^{-1/N(0)V}$ and all its
derivatives vanish as $V\to0^+$. What does this imply about
perturbation theory in $V$?
\begin{solution}
Let $x=N(0)V\to0^+$. Then $\Delta\propto e^{-1/x}$. Every derivative
$\dd^n\Delta/\dd x^n$ is a polynomial in $1/x$ times $e^{-1/x}$, and
$e^{-1/x}$ decays faster than any power of $x$, so all derivatives
$\to0$ as $x\to0^+$. The Taylor series of $\Delta$ about $V=0$ is
therefore identically zero, yet $\Delta\neq0$ for $V>0$: $\Delta$ is
non-analytic at $V=0$. No finite order of perturbation theory in $V$
can produce the gap --- the superconducting state is invisible to an
expansion around the normal metal, which is why a non-perturbative
ansatz was essential.
\end{solution}
\end{exercise}

\begin{exercise}[The gap as a detector threshold]
Aluminium has $2\Delta\approx0.34\,$meV. What is the longest-wavelength
photon that can break a Cooper pair? Compare to the $5\,$GHz operating
frequency of a typical transmon and comment on quasiparticle
poisoning.
\begin{solution}
A photon breaks a pair if $\hbar\omega\geq2\Delta$, i.e.\
$\omega\geq2\Delta/\hbar$. With $2\Delta=0.34\,$meV,
$\nu=2\Delta/h=0.34\times10^{-3}\,\text{eV}/(4.14\times10^{-15}\,\text{eV·s})
\approx82\,$GHz, or wavelength $c/\nu\approx3.6\,$mm
(sub-THz/far-infrared). A $5\,$GHz transmon drive is a factor
$\sim16$ below this threshold, so it cannot directly break pairs ---
good. But stray high-frequency radiation, cosmic rays, or
pair-breaking phonons above $82\,$GHz inject quasiparticles that
tunnel across junctions and relax the qubit (\cref{ch:qp-detection}):
the same gap that protects the qubit at its operating frequency sets
the threshold above which the environment poisons it.
\end{solution}
\end{exercise}

\begin{exercise}[$LC$ equation of motion from the Lagrangian]
Derive the Euler--Lagrange equation of the $LC$ Lagrangian
\cref{eq:lc-lagrangian} and confirm it gives $\omega=1/\sqrt{LC}$.
Identify the mechanical analogue of each term.
\begin{solution}
$\partial\mathcal L/\partial\dot\Phi=C\dot\Phi$,
$\partial\mathcal L/\partial\Phi=-\Phi/L$. Euler--Lagrange:
$\frac{\dd}{\dd t}(C\dot\Phi)-(-\Phi/L)=C\ddot\Phi+\Phi/L=0$, i.e.\
$\ddot\Phi+(1/LC)\Phi=0$, harmonic motion at
$\omega=1/\sqrt{LC}$. The term $\tfrac12 C\dot\Phi^2$ is kinetic
energy ($C\leftrightarrow$ mass, $\dot\Phi=V\leftrightarrow$ velocity)
and $\Phi^2/2L$ is potential energy
($1/L\leftrightarrow$ spring constant, $\Phi\leftrightarrow$
displacement).
\end{solution}
\end{exercise}

\begin{exercise}[Zero-point fluctuations and impedance]
Starting from the oscillator analogy ($C\leftrightarrow m$,
$\omega=1/\sqrt{LC}$), derive $\Phi_{\rm zpf}=\sqrt{\hbar Z/2}$ and
$Q_{\rm zpf}=\sqrt{\hbar/2Z}$ with $Z=\sqrt{L/C}$, and verify
$\Phi_{\rm zpf}Q_{\rm zpf}=\hbar/2$.
\begin{solution}
For a mechanical oscillator $x_{\rm zpf}=\sqrt{\hbar/2m\omega}$. With
$m\to C$, $\omega\to1/\sqrt{LC}$:
$\Phi_{\rm zpf}=\sqrt{\hbar/(2C\cdot1/\sqrt{LC})}
=\sqrt{\hbar\sqrt{LC}/2C}=\sqrt{(\hbar/2)\sqrt{L/C}}=\sqrt{\hbar Z/2}$.
Likewise $p_{\rm zpf}=\sqrt{\hbar m\omega/2}$ gives
$Q_{\rm zpf}=\sqrt{\hbar C/(2\sqrt{LC})}=\sqrt{(\hbar/2)\sqrt{C/L}}
=\sqrt{\hbar/2Z}$. Their product is
$\sqrt{(\hbar Z/2)(\hbar/2Z)}=\sqrt{\hbar^2/4}=\hbar/2$, saturating the
uncertainty bound as the vacuum of a harmonic oscillator must.
\end{solution}
\end{exercise}

\begin{exercise}[Photon energy versus temperature]
For an $LC$ resonator at $\omega/2\pi=5\,$GHz, compute the
single-photon energy in $\mu$eV and the equivalent temperature
$\hbar\omega/k_B$. Why must the experiment run in a dilution
refrigerator?
\begin{solution}
$\hbar\omega=h\nu=(4.14\times10^{-15}\,\text{eV·s})(5\times10^9\,\text{Hz})
\approx2.07\times10^{-5}\,\text{eV}=20.7\,\mu$eV. The temperature is
$\hbar\omega/k_B=h\nu/k_B=(5\times10^9)/(2.08\times10^{10}\,\text{Hz/K})
\approx0.24\,$K$=240\,$mK. To keep the resonator in its ground state
the thermal occupation $\bar n=1/(e^{\hbar\omega/k_BT}-1)$ must be
$\ll1$, which needs $T\ll240\,$mK; a dilution refrigerator at
$10$--$20\,$mK gives $\bar n\sim10^{-5}$, justifying the vacuum-state
assumption of Part~III.
\end{solution}
\end{exercise}

\begin{exercise}[Why the harmonic ladder forbids a qubit]
Explain quantitatively why a resonant drive on a linear $LC$ resonator
cannot prepare an arbitrary superposition confined to $\{\ket0,\ket1\}$,
and state what property of the spectrum must change to allow it.
\begin{solution}
The spectrum is $E_n=\hbar\omega(n+\tfrac12)$, so every neighbouring
transition has the same frequency $\omega$: $E_1-E_0=E_2-E_1=\cdots
=\hbar\omega$. A drive at $\omega$ resonantly couples $\ket0\!\leftrightarrow\!\ket1$
but \emph{equally} couples $\ket1\!\leftrightarrow\!\ket2$ and higher, so
population leaks out of the two-level subspace --- one cannot perform an
independent rotation on $\{\ket0,\ket1\}$. To isolate the lowest
transition the spectrum must be \emph{anharmonic}: the
$\ket1\to\ket2$ frequency must differ from $\ket0\to\ket1$ by more than
the drive bandwidth. The anharmonicity $\alpha=(E_2-E_1)-(E_1-E_0)$
must be nonzero, which a linear element cannot supply --- hence the
Josephson junction of \cref{ch:josephson}.
\end{solution}
\end{exercise}
